% Part: proof-theory
% Chapter: natural-deduction
% Section: quantifiers

\documentclass[../../../include/open-logic-section]{subfiles}

\begin{document}

\olfileid[id]{pt}{nat}{qua}

\olsection{\usetoken{P}{proof} Reguler dan Substitusi}

Aturan-aturan kuantor menimbulkan sejumlah kerumitan karena
``syarat variabel eigen'' yang berlaku pada aturan
\Elim{\lexists} dan~\Intro{\lforall}. Sebagian besar kerumitan ini
dapat ditangani dengan memastikan bahwa !!{constant}~$c$ dalam
inferensi-inferensi tersebut tidak pernah digunakan kembali. Mari kita
perkenalkan definisi berikut.

\begin{defn}
Suatu !!{proof} deduksi alami~$\delta$ disebut \emph{reguler} jika
\begin{enumerate}
  \item tidak ada variabel eigen~$a$ yang digunakan sebagai variabel eigen
  bagi lebih dari satu inferensi \Elim{\lexists}
  atau~\Intro{\lforall}; dan
  \item jika $a$ adalah variabel eigen bagi suatu inferensi
  \Elim{\lexists} atau~\Intro{\lforall}, maka $a$ hanya muncul dalam
  sub-!!{proof} langsung yang masing-masing menuju premis minor atau
  premis inferensi tersebut.
\end{enumerate}
\end{defn}

\begin{lem}\ollabel{lem:subst-regular}
Jika $\delta$ reguler dan berakhir pada~$!C$, $c$ adalah
!!a{constant} yang tidak digunakan sebagai variabel eigen
dalam~$\delta$, dan $t$ adalah suku yang tidak memuat variabel eigen
apa pun dari~$\delta$, maka $\Subst{\delta}{t}{c}$ yang diperoleh
dari~$\delta$ dengan mengganti setiap kemunculan~$c$ dengan~$t$
merupakan !!{proof} reguler. !!{formula} akhir
$\Subst{\delta}{t}{c}$ ialah $\Subst{!C}{t}{c}$. Jika $!A^x$
merupakan asumsi terbuka dari~$\delta$, maka
$\Subst{!A}{t}{c}^x$ merupakan asumsi terbuka dari
$\Subst{\delta}{t}{c}$.
\end{lem}

\begin{proof}
  Melalui induksi pada $\pheight{\delta}$. Jika
  $\pheight{\delta}=0$, $\delta$ hanya terdiri atas asumsi~$!A^x$.
  Maka $\Subst{\delta}{t}{c}$ ialah~$\Subst{!A}{t}{c}^x$.

  Jika $\pheight{\delta}>0$, kita membedakan kasus menurut inferensi
  terakhir dalam~$\delta$. Kasus-kasus aturan penghubung proposisional
  mudah dan ditinggalkan sebagai latihan. Kita membahas kasus ketika
  $\delta$ berakhir pada inferensi kuantor.
  \begin{enumerate}
    \item Jika inferensi terakhir adalah $\Intro{\lforall}$, $\delta$
    berbentuk
    \[
    \AxiomC{}
    \RightLabel{$\delta'$}
    \DeduceC{$!A(a,c)$}
    \RightLabel{\Intro{\lforall}}
    \UnaryInfC{$\lforall[x][!A(x,c)]$}
    \DisplayProof
    \]
    Berdasarkan hipotesis induksi, $\Subst{\delta'}{t}{c}$ merupakan
    !!{proof} reguler atas $!A(a,t)$. Karena $a$ merupakan variabel
    eigen dari~$\delta$, $a$ tidak muncul dalam~$t$. Karena itu,
    inferensi terakhir dalam $\Subst{\delta}{t}{c}$,
    \[
    \AxiomC{}
    \RightLabel{$\Subst{\delta'}{t}{c}$}
    \DeduceC{$!A(a,t)$}
    \RightLabel{\Intro{\lforall}}
    \UnaryInfC{$\lforall[x][!A(x,t)]$}
    \DisplayProof
    \]
    adalah benar: karena $t$ tidak memuat~$a$, baik kesimpulan maupun
    asumsi terbuka mana pun tidak memuat~$a$.
    \item Jika inferensi terakhir adalah $\Elim{\lforall}$, $\delta$
    berbentuk
    \[
    \AxiomC{}
    \RightLabel{$\delta'$}
    \DeduceC{$\lforall[x][!A(x,c)]$}
    \RightLabel{\Elim{\lforall}}
    \UnaryInfC{$!A(s(c),c)$}
    \DisplayProof
    \]
    Berdasarkan hipotesis induksi, $\Subst{\delta'}{t}{c}$ merupakan
    !!{proof} reguler atas $\lforall[x][!A(x,t)]$. (Suku~$s(c)$
    dalam kesimpulan boleh memuat~$c$, tetapi tidak harus.) Inferensi
    terakhir dalam $\Subst{\delta}{t}{c}$,
    \[
    \AxiomC{}
    \RightLabel{$\Subst{\delta'}{t}{c}$}
    \DeduceC{$\lforall[x][!A(x,t)]$}
    \RightLabel{\Elim{\lforall}}
    \UnaryInfC{$!A(s(t),t)$}
    \DisplayProof
    \]
    adalah benar, dan
    $!A(s(t),t) \ident \Subst{!A(s(c),c)}{t}{c}$. Kasus
    \Intro{\lexists} serupa.
    \item Jika inferensi terakhir adalah $\Elim{\lexists}$, $\delta$
    berbentuk
    \[
    \AxiomC{}
    \RightLabel{$\delta_1'$}
    \DeduceC{$\lexists[x][!A(x,c)]$}
    \AxiomC{$\Discharge{!A(a,c)}{x}$}
    \RightLabel{$\delta_2'$}
    \DeduceC{$!C$}
    \DischargeRule{\Elim{\lexists}}{x}
    \BinaryInfC{$!C$}
    \DisplayProof
    \]
    Berdasarkan hipotesis induksi,
    $\Subst{\delta_1'}{t}{c}$ dan
    $\Subst{\delta_2'}{t}{c}$ masing-masing merupakan !!{proof}
    reguler atas $\lexists[x][!A(x,t)]$ dan atas
    $\Subst{!C}{t}{c}$ dari asumsi terbuka~$!A(a,t)$. Karena $a$
    merupakan variabel eigen dari~$\delta$, $a$ tidak muncul
    dalam~$t$. Maka inferensi terakhir dalam
    $\Subst{\delta}{t}{c}$,
    \[
    \AxiomC{}
    \RightLabel{$\Subst{\delta_1'}{t}{c}$}
    \DeduceC{$\lexists[x][!A(x,t)]$}
    \AxiomC{$\Discharge{!A(a,t)}{x}$}
    \RightLabel{$\Subst{\delta_2'}{t}{c}$}
    \DeduceC{$\Subst{!C}{t}{c}$}
    \DischargeRule{\Elim{\lexists}}{x}
    \BinaryInfC{$\Subst{!C}{t}{c}$}
    \DisplayProof
    \]
    adalah benar:
    $\Subst{!A(x,t)}{a}{x}
    \ident \Subst{!A(a,c)}{t}{c}$, dan baik kesimpulan
    $\Subst{!C}{t}{c}$ maupun asumsi terbuka mana pun tidak
    memuat~$a$.
  \end{enumerate}
\end{proof}

\begin{prop}\ollabel{prop:regular}
Jika $\delta$ adalah !!a{proof} dalam \Log{N1c} atau \Log{N1i}, maka
terdapat !!{proof} reguler~$\delta'$ dengan struktur yang sama (khususnya,
dengan !!{height} yang sama) yang berbeda dari $\delta$ hanya melalui
penggantian variabel-variabel eigen.
\end{prop}

\begin{proof}
Khusus untuk keperluan bukti ini, sebut suatu inferensi variabel eigen
dalam~$\delta$ \emph{bersih} jika variabel eigennya bukan variabel eigen
bagi inferensi lain dan tidak muncul di tempat mana pun dalam~$\delta$
selain sub-!!{proof} langsung dari inferensi tersebut; sebut inferensi itu
\emph{kotor} jika tidak demikian. Misalkan $n$ banyaknya inferensi
variabel eigen kotor dalam~$\delta$. Kita menunjukkan melalui induksi
pada~$n$ bahwa $\delta$ dapat diubah menjadi !!{proof} reguler~$\delta'$
yang disyaratkan. Jika $n=0$, semua inferensi variabel eigen
dalam~$\delta$ bersih; jadi $\delta$ reguler dan tidak ada yang perlu
dibuktikan. Jika tidak, pilih inferensi variabel eigen kotor paling atas,
misalnya \Intro{\lforall}. Sub-!!{proof}~$\delta_1$ dari~$\delta$ yang
berakhir pada inferensi itu berbentuk
    \[
    \AxiomC{}
    \RightLabel{$\delta_2$}
    \DeduceC{$!A(c)$}
    \RightLabel{\Intro{\lforall}}
    \UnaryInfC{$\lforall[x][!A(x)]$}
    \DisplayProof
    \]
Karena $c$ merupakan variabel eigen, $c$ tidak muncul
dalam~$\lforall[x][!A(x)]$ ataupun dalam asumsi terbuka. Tidak ada
inferensi variabel eigen kotor dalam $\delta_2$, sehingga $\delta_2$
reguler; khususnya, $c$ tidak digunakan sebagai variabel eigen di
dalamnya. Misalkan $a$ adalah !!a{constant} yang tidak muncul
dalam~$\delta$. Berdasarkan \olref{lem:subst-regular},
$\Subst{\delta_2}{a}{c}$ merupakan bukti reguler atas $!A(a)$.
Berdasarkan syarat variabel eigen, tidak ada asumsi terbuka
dalam~$\delta_2$ yang memuat~$c$; karena itu, tidak ada asumsi terbuka
dalam~$\Subst{\delta_2}{a}{c}$ yang memuat~$a$. Jadi kita dapat
menerapkan~$\Intro{\lforall}$. Jika kita mengganti $\delta_1$
dalam~$\delta$ dengan bukti $\delta_1'$,
    \[
    \AxiomC{}
    \RightLabel{$\Subst{\delta_2}{a}{c}$}
    \DeduceC{$!A(a)$}
    \RightLabel{\Intro{\lforall}}
    \UnaryInfC{$\lforall[x][!A(x)]$}
    \DisplayProof
    \]
kita telah mengganti satu inferensi variabel eigen kotor dengan inferensi
bersih, dan satu-satunya perbedaan ialah penggantian variabel eigen $c$
dengan~$a$. Bukti yang dihasilkan memiliki lebih sedikit daripada $n$
inferensi variabel eigen kotor, sehingga hasilnya mengikuti dari hipotesis
induksi.
\end{proof}

\begin{prob}
Uraikan kasus dalam bukti \olref{prop:regular} ketika inferensi variabel
eigen paling atas adalah~\Elim{\lexists}.
\end{prob}

Kita tidak akan merujuk secara eksplisit pada proposisi yang baru saja
dibuktikan, tetapi kita akan mengasumsikan bahwa semua !!{proof} yang
dibahas adalah reguler---jika tidak, kita selalu dapat mengganti variabel
eigen untuk menjadikannya reguler.

\Olref{lem:subst-regular} dan \olref{prop:regular} juga berlaku bagi
\Log{N2c} dan~\Log{N2i}. Bukti-buktinya ditinggalkan sebagai latihan.

\end{document}
